\documentclass[10pt,a4paper]{amsart}
\usepackage{amsmath,amssymb,amsthm,graphicx}
\usepackage{verbatim}
\newenvironment{reference}{\comment}{\endcomment}
%\newenvironment{reference}{}{}
\newenvironment{slogan}{\comment}{\endcomment}
\newenvironment{history}{\comment}{\endcomment}

% For commutative diagrams we use Xy-pic
\usepackage[all]{xy}

% We use 2cell for 2-commutative diagrams.
\xyoption{2cell}
% Theorem environments.
%
\theoremstyle{plain}
\newtheorem{theorem}[subsection]{Theorem}
\newtheorem{proposition}[subsection]{Proposition}
\newtheorem{lemma}[subsection]{Lemma}

\theoremstyle{definition}
\newtheorem{definition}[subsection]{Definition}
\newtheorem{example}[subsection]{Example}
\newtheorem{exercise}[subsection]{Exercise}
\newtheorem{situation}[subsection]{Situation}

\theoremstyle{remark}
\newtheorem{remark}[subsection]{Remark}
\newtheorem{remarks}[subsection]{Remarks}

\numberwithin{equation}{subsection}

% Macros
%
\def\lim{\mathop{\mathrm{lim}}\nolimits}
\def\colim{\mathop{\mathrm{colim}}\nolimits}
\def\Spec{\mathop{\mathrm{Spec}}}
\def\Hom{\mathop{\mathrm{Hom}}\nolimits}
\def\Ext{\mathop{\mathrm{Ext}}\nolimits}
\def\SheafHom{\mathop{\mathcal{H}\!\mathit{om}}\nolimits}
\def\SheafExt{\mathop{\mathcal{E}\!\mathit{xt}}\nolimits}
\def\Sch{\mathit{Sch}}
\def\Mor{\mathop{\mathrm{Mor}}\nolimits}
\def\Ob{\mathop{\mathrm{Ob}}\nolimits}
\def\Sh{\mathop{\mathit{Sh}}\nolimits}
\def\NL{\mathop{N\!L}\nolimits}
\def\CH{\mathop{\mathrm{CH}}\nolimits}
\def\proetale{{pro\text{-}\acute{e}tale}}
\def\etale{{\acute{e}tale}}
\def\QCoh{\mathit{QCoh}}
\def\Ker{\mathop{\mathrm{Ker}}}
\def\Im{\mathop{\mathrm{Im}}}
\def\Coker{\mathop{\mathrm{Coker}}}
\def\Coim{\mathop{\mathrm{Coim}}}

% Boxtimes
%
\DeclareMathSymbol{\boxtimes}{\mathbin}{AMSa}{"02}

%
% Macros for moduli stacks/spaces
%
\def\QCohstack{\mathcal{QC}\!\mathit{oh}}
\def\Cohstack{\mathcal{C}\!\mathit{oh}}
\def\Spacesstack{\mathcal{S}\!\mathit{paces}}
\def\Quotfunctor{\mathrm{Quot}}
\def\Hilbfunctor{\mathrm{Hilb}}
\def\Curvesstack{\mathcal{C}\!\mathit{urves}}
\def\Polarizedstack{\mathcal{P}\!\mathit{olarized}}
\def\Complexesstack{\mathcal{C}\!\mathit{omplexes}}
% \Pic is the operator that assigns to X its picard group, usage \Pic(X)
% \Picardstack_{X/B} denotes the Picard stack of X over B
% \Picardfunctor_{X/B} denotes the Picard functor of X over B
\def\Pic{\mathop{\mathrm{Pic}}\nolimits}
\def\Picardstack{\mathcal{P}\!\mathit{ic}}
\def\Picardfunctor{\mathrm{Pic}}
\def\Deformationcategory{\mathcal{D}\!\mathit{ef}}
\usepackage{cleveref}
\crefname{originallemma}{original lemma}{original lemmas}
\Crefname{originallemma}{Original lemma}{Original lemmas}
\makeatletter
\newlabel{morphisms-lemma-finite-fibre}{{02NG}{source}{}{sourcepointer-1}{}}
\newlabel{morphisms-lemma-finite-fibre@cref}{{[originallemma][0][]02NG}{[1][0][]source}{}{}{}}
\makeatother
\crefname{originallemma}{original lemma}{original lemmas}
\Crefname{originallemma}{Original lemma}{Original lemmas}
\makeatletter
\newlabel{morphisms-lemma-permanence-finite-type}{{01T8}{source}{}{sourcepointer-2}{}}
\newlabel{morphisms-lemma-permanence-finite-type@cref}{{[originallemma][0][]01T8}{[1][0][]source}{}{}{}}
\makeatother
\crefname{originallemma}{original lemma}{original lemmas}
\Crefname{originallemma}{Original lemma}{Original lemmas}
\makeatletter
\newlabel{schemes-lemma-monomorphism-separated}{{01L4}{source}{}{sourcepointer-3}{}}
\newlabel{schemes-lemma-monomorphism-separated@cref}{{[originallemma][0][]01L4}{[1][0][]source}{}{}{}}
\makeatother
\crefname{originalsection}{original section}{original sections}
\Crefname{originalsection}{Original section}{Original sections}
\makeatletter
\newlabel{section-representable}{{03HA}{source}{}{sourcepointer-4}{}}
\newlabel{section-representable@cref}{{[originalsection][0][]03HA}{[1][0][]source}{}{}{}}
\makeatother
\crefname{originallemma}{original lemma}{original lemmas}
\Crefname{originallemma}{Original lemma}{Original lemmas}
\makeatletter
\newlabel{spaces-lemma-properties-diagonal}{{02X4}{source}{}{sourcepointer-5}{}}
\newlabel{spaces-lemma-properties-diagonal@cref}{{[originallemma][0][]02X4}{[1][0][]source}{}{}{}}
\makeatother
\crefname{originallemma}{original lemma}{original lemmas}
\Crefname{originallemma}{Original lemma}{Original lemmas}
\makeatletter
\newlabel{spaces-lemma-representable-morphisms-spaces-property}{{02WZ}{source}{}{sourcepointer-6}{}}
\newlabel{spaces-lemma-representable-morphisms-spaces-property@cref}{{[originallemma][0][]02WZ}{[1][0][]source}{}{}{}}
\makeatother
\title{Five properties of the diagonal morphism of algebraic spaces}
\author{The Stacks Project Authors}
\date{stacks-project:spaces-morphisms.tex:local-snapshot-20261012; GNU FDL 1.2 or later}
\begin{document}
\maketitle
\setcounter{subsection}{1}
\noindent\textit{Editorial note.} All delivered mathematical text is the author's own, copied from the pinned Stacks Project chapter file spaces-morphisms.tex; no statement or proof step was rewritten, reordered or removed.\par\medskip
\noindent\textit{Editorial note.} The proof recalls a computation performed in the proof of a cross-chapter lemma and then carries out the transfer for the relative diagonal explicitly, property by property.\par\medskip
\noindent\textit{Editorial note.} The six references of the counted spans are recorded as non-load-bearing navigation pointers; the argument that consumes them stays inside the counted proof span.\par\medskip
\noindent\textit{Editorial note.} Preamble: the three native excerpt spans of sources/stacks-spaces-morphisms/preamble.tex (lines 8-13, 14-18 and 154-217). The chapter file itself carries no preamble text.\par\medskip
\noindent\textit{Allowed standard prerequisites (editorial).}\begin{itemize}
\item The diagonal of an algebraic space over its base is representable with properties (2)-(5). Spaces, Lemma spaces-lemma-properties-diagonal, used at proof lines 397-400 and 432-434. Supplies the base case that the whole argument transfers.
\item Properties of morphisms of algebraic spaces which are etale local on the source-and-target. Spaces, Lemma spaces-lemma-representable-morphisms-spaces-property, used at proof lines 429-432. Reduces the five claims to the corresponding claims for R -> U x\_Y U.
\item A composition of monomorphisms is a monomorphism and a monomorphism is separated. Schemes, Lemma schemes-lemma-monomorphism-separated, used at proof lines 440-442 and 449-450. Gives properties (3) and (4).
\item Permanence of local finite typeness under composition. Morphisms, Lemma morphisms-lemma-permanence-finite-type, used at proof lines 443-445. Gives property (2) for R -> U x\_Y U.
\item A monomorphism locally of finite type has finite fibres, hence is locally quasi-finite. Morphisms, Lemma morphisms-lemma-finite-fibre, used at proof lines 446-448. Gives property (5).
\end{itemize}\medskip
\ifdefined\hypertarget\hypertarget{sourcepointer-1}{}\fi
\noindent\textit{Source reference.} Original lemma 02NG: focused-seeds/AGP/sources/stacks-snapshot-20261011/tags.txt, source lines 3479--3479. Navigation pointer into the Stacks Project; the standard fact it supplies is applied inside the counted proof span and the argument resting on it stays there.\par\medskip
\ifdefined\hypertarget\hypertarget{sourcepointer-2}{}\fi
\noindent\textit{Source reference.} Original lemma 01T8: focused-seeds/AGP/sources/stacks-snapshot-20261011/tags.txt, source lines 2365--2365. Navigation pointer into the Stacks Project; the standard fact it supplies is applied inside the counted proof span and the argument resting on it stays there.\par\medskip
\ifdefined\hypertarget\hypertarget{sourcepointer-3}{}\fi
\noindent\textit{Source reference.} Original lemma 01L4: focused-seeds/AGP/sources/stacks-snapshot-20261011/tags.txt, source lines 2073--2073. Navigation pointer into the Stacks Project; the standard fact it supplies is applied inside the counted proof span and the argument resting on it stays there.\par\medskip
\ifdefined\hypertarget\hypertarget{sourcepointer-4}{}\fi
\noindent\textit{Source reference.} Original section 03HA: focused-seeds/AGP/sources/stacks-snapshot-20261011/tags.txt, source lines 4559--4559. Navigation pointer into the Stacks Project; the standard fact it supplies is applied inside the counted proof span and the argument resting on it stays there.\par\medskip
\ifdefined\hypertarget\hypertarget{sourcepointer-5}{}\fi
\noindent\textit{Source reference.} Original lemma 02X4: focused-seeds/AGP/sources/stacks-snapshot-20261011/tags.txt, source lines 3827--3827. Navigation pointer into the Stacks Project; the standard fact it supplies is applied inside the counted proof span and the argument resting on it stays there.\par\medskip
\ifdefined\hypertarget\hypertarget{sourcepointer-6}{}\fi
\noindent\textit{Source reference.} Original lemma 02WZ: focused-seeds/AGP/sources/stacks-snapshot-20261011/tags.txt, source lines 3822--3822. Navigation pointer into the Stacks Project; the standard fact it supplies is applied inside the counted proof span and the argument resting on it stays there.\par\medskip
\begin{lemma}
\label{lemma-properties-diagonal}
Let $S$ be a scheme contained in $\Sch_{fppf}$.
Let $f : X \to Y$ be a morphism of algebraic spaces over $S$.
Let $\Delta_{X/Y} : X \to X \times_Y X$ be the diagonal morphism.
Then
\begin{enumerate}
\item $\Delta_{X/Y}$ is representable,
\item $\Delta_{X/Y}$ is locally of finite type,
\item $\Delta_{X/Y}$ is a monomorphism,
\item $\Delta_{X/Y}$ is separated, and
\item $\Delta_{X/Y}$ is locally quasi-finite.
\end{enumerate}
\end{lemma}

\begin{proof}
We are going to use the fact that $\Delta_{X/S}$ is
representable (by definition of an algebraic space) and that
it satisfies properties (2) -- (5), see
Spaces, Lemma \ref{spaces-lemma-properties-diagonal}.
Note that we have a factorization
$$
X
\longrightarrow
X \times_Y X
\longrightarrow
X \times_S X
$$
of the diagonal $\Delta_{X/S} : X \to X \times_S X$. Since
$X \times_Y X \to X \times_S X$ is a monomorphism, and since
$\Delta_{X/S}$ is representable, it follows formally that
$\Delta_{X/Y}$ is representable. In particular, the rest of
the statements now make sense, see
Section \ref{section-representable}.

\medskip\noindent
Choose a surjective \'etale morphism $U \to X$, with $U$ a scheme.
Consider the diagram
$$
\xymatrix{
R = U \times_X U \ar[r] \ar[d] &
U \times_Y U \ar[d] \ar[r] &
U \times_S U \ar[d] \\
X \ar[r] & X \times_Y X \ar[r] & X \times_S X
}
$$
Both squares are cartesian, hence so is the outer rectangle.
The top row consists of schemes, and the vertical arrows
are surjective \'etale morphisms. By
Spaces, Lemma \ref{spaces-lemma-representable-morphisms-spaces-property}
the properties (2) -- (5) for $\Delta_{X/Y}$ are equivalent to those of
$R \to U \times_Y U$. In the proof of
Spaces, Lemma \ref{spaces-lemma-properties-diagonal}
we have seen that $R \to U \times_S U$ has properties (2) -- (5).
The morphism $U \times_Y U \to U \times_S U$ is a monomorphism
of schemes. These facts imply that $R \to U \times_Y U$ have
properties (2) -- (5).

\medskip\noindent
Namely: For (3), note that $R \to U \times_Y U$
is a monomorphism as the composition
$R \to U \times_S U$ is a monomorphism. For (2), note that
$R \to U \times_Y U$ is locally of finite type, as the
composition $R \to U \times_S U$ is locally of finite type
(Morphisms, Lemma \ref{morphisms-lemma-permanence-finite-type}).
A monomorphism which is locally of finite type is locally quasi-finite
because it has finite fibres
(Morphisms, Lemma \ref{morphisms-lemma-finite-fibre}), hence (5).
A monomorphism is separated
(Schemes, Lemma \ref{schemes-lemma-monomorphism-separated}), hence (4).
\end{proof}

\end{document}
